% El documento usa subíndices y símbolos Unicode en la prosa (PM₂.₅, NO₂, O₃, ≥, −, →). La fuente
% por defecto de LaTeX no los tiene y saldrían como huecos: "PM ." en vez de "PM₂.₅". Cada uno se
% mapea a su equivalente de LaTeX para que el PDF diga lo mismo que el HTML.
\usepackage{newunicodechar}
\newunicodechar{₂}{\textsubscript{2}}
\newunicodechar{₃}{\textsubscript{3}}
\newunicodechar{₅}{\textsubscript{5}}
\newunicodechar{⁵}{\textsuperscript{5}}
\newunicodechar{−}{\ensuremath{-}}
\newunicodechar{≥}{\ensuremath{\geq}}
\newunicodechar{≤}{\ensuremath{\leq}}
\newunicodechar{≈}{\ensuremath{\approx}}
\newunicodechar{→}{\ensuremath{\rightarrow}}
\newunicodechar{√}{\ensuremath{\surd}}
\newunicodechar{Σ}{\ensuremath{\Sigma}}
\newunicodechar{α}{\ensuremath{\alpha}}
\newunicodechar{β}{\ensuremath{\beta}}
\newunicodechar{γ}{\ensuremath{\gamma}}
\newunicodechar{ν}{\ensuremath{\nu}}
\newunicodechar{ρ}{\ensuremath{\rho}}
\newunicodechar{ŷ}{\ensuremath{\hat{y}}}
\newunicodechar{…}{\ldots}
